%% ============================================================
%%  HQDE — IEEE Transactions Journal Paper  (v3 — 20 pages)
%%  Compile:  pdflatex -> pdflatex -> pdflatex
%% ============================================================
\documentclass[journal,twoside]{IEEEtran}

%% ---- Packages -----------------------------------------------
\usepackage[T1]{fontenc}
\usepackage[utf8]{inputenc}
\usepackage{xspace}
\usepackage{amsmath,amssymb,amsfonts}
\usepackage{booktabs}
\usepackage{multirow}
\usepackage{array}
\usepackage{tabularx}
\usepackage{adjustbox}
\usepackage{graphicx}
\usepackage{xcolor}
\usepackage{url}
\usepackage{cite}
\usepackage{microtype}
\usepackage{balance}
\usepackage{subcaption}
\usepackage[ruled,vlined,linesnumbered]{algorithm2e}
\usepackage[hidelinks]{hyperref}

%% ---- Custom column types ------------------------------------
\newcolumntype{L}[1]{>{\raggedright\arraybackslash}p{#1}}
\newcolumntype{C}[1]{>{\centering\arraybackslash}p{#1}}
\newcolumntype{R}[1]{>{\raggedleft\arraybackslash}p{#1}}

%% ---- Macros -------------------------------------------------
\newcommand{\hqde}{\mbox{HQDE}\xspace}
\newcommand{\eg}{\textit{e.g.},\xspace}
\newcommand{\ie}{\textit{i.e.},\xspace}
\newcommand{\etal}{\textit{et al.}\xspace}
\DeclareMathOperator*{\argmax}{arg\,max}

%% ============================================================
\begin{document}

\title{HQDE: Hierarchical Quantized Distributed Ensemble Learning\\
       with Adaptive Communication Compression}

\author{Prathamesh~Nikam and P.S.V.S.~Sai~Prasad%
\thanks{Both authors are with the School of Computer and Information
Sciences, University of Hyderabad, Hyderabad~500046, India.
Corresponding author: Prathamesh Nikam
(e-mail: prathamesh.nikam@uohyd.ac.in).}%
\thanks{Manuscript received 2025.}}

\markboth{IEEE Transactions on Neural Networks and Learning Systems}%
{Nikam \& Sai~Prasad: HQDE}

\maketitle

%% ============================================================
%%  ABSTRACT
%% ============================================================

\begin{abstract}
Distributed ensemble training must balance model diversity, communication cost, and stability under heterogeneous data. Existing frameworks often handle these goals in isolation. We implement \hqde{}, a distributed ensemble system that trains multiple workers in parallel, tracks per-worker efficiency, and optionally performs epoch-level federated averaging with adaptive delta quantization. The implementation uses Ray actors when Ray is available and falls back to local workers otherwise, so each worker keeps its model and optimizer state in memory. First, workers run local mini-batch updates and update an efficiency score based on accuracy and loss. Next, in federated mode, workers send weight deltas that are compressed with blockwise symmetric quantization, scheduled bitwidths, and error feedback after a warmup period. Then, the server aggregates deltas with a weighted mean or a FedAdam-style update and broadcasts the new reference weights. Finally, inference uses efficiency-weighted softmax or mean fusion of logits. We report CIFAR-10 example results from the provided scripts and discuss trade-offs across aggregation, quantization, and batch-normalization policies.
\end{abstract}

\begin{IEEEkeywords}
distributed deep learning, ensemble methods, federated learning,
communication compression, adaptive quantization, Ray framework,
stateful actors, batch normalization, non-IID data.
\end{IEEEkeywords}

%% ============================================================
%%  I. INTRODUCTION
%% ============================================================
\section{Introduction}
\label{sec:intro}

The scale of modern deep learning has made distributed training
indispensable. A single ResNet-50 requires approximately $4\times10^9$
floating-point operations per image and several days on a single GPU
to converge on ImageNet; scaling this to real-world datasets with
hundreds of millions of samples is only tractable through parallelism.
Yet distributing training introduces a set of structural challenges
that raw computational power alone cannot resolve.

\subsection{The Three-Way Problem}

Three requirements must be met simultaneously for effective distributed
deep learning.

\textbf{Accuracy through diversity.} A single model, no matter how
large, is constrained by its initialization point and training
trajectory. Research in ensemble learning~\cite{lakshminarayanan2017}
consistently shows that combining predictions from multiple
independently trained models improves accuracy and calibration beyond
what any individual model can achieve. The mechanism is model diversity:
networks initialized differently and optimized along different
trajectories develop complementary representations and make different
errors. When their predictions are combined, errors that are
idiosyncratic to any one model cancel out. This is the core principle
behind Random Forests~\cite{breiman2001}, Deep
Ensembles~\cite{lakshminarayanan2017}, and Snapshot
Ensembles~\cite{huang2017snapshot}, and it is the primary motivation
for the ensemble-based design of \hqde{}.

\textbf{Communication efficiency.} Each synchronization round in a
distributed system requires transmitting model weights between workers
and a central coordinator. For a ResNet-18 with $11.17\times10^6$
parameters at 32-bit precision, each transmission requires approximately
44~MB. Over 20 training rounds with 4 workers, the total communication
volume exceeds 4.4~GB. In bandwidth-limited settings such as mobile
networks or cross-institutional collaborations, this volume is
prohibitive. Communication compression techniques reduce this cost, but
existing methods~\cite{alistarh2017,bernstein2018} were designed as
standalone gradient modifiers and do not integrate cleanly with ensemble
training or federated aggregation.

\textbf{Robustness to data heterogeneity.} In federated and
edge-computing settings, training data is distributed across nodes that
may have very different class distributions. A medical imaging node may
have abundant chest X-rays but no brain scans; a mobile keyboard node
may have typing data from users who primarily speak one language. This
\textbf{non-IID data distribution} causes client models to optimize
for local objectives that diverge from the global objective, a
phenomenon known as \textbf{client drift}~\cite{karimireddy2020scaffold}.
Standard aggregation strategies that assume uniform data distributions
degrade significantly under these conditions.

\subsection{Limitations of Existing Approaches}

To the best of our knowledge, no single framework simultaneously addresses all three requirements.

\textbf{Ensemble methods}~\cite{lakshminarayanan2017,huang2017snapshot}
maximize accuracy through model diversity but make no attempt to reduce
communication overhead. Every member trains on the full dataset,
requiring full dataset replication across all workers. Privacy
preservation is not a design goal.

\textbf{Federated learning algorithms}~\cite{mcmahan2017,li2020fedprox,
karimireddy2020scaffold} enable training without data centralization and
were specifically designed for non-IID settings. However, they converge
to a single global model shared across all clients, which cannot provide
the accuracy benefits of ensembling. Per-round communication volume
equals the full model weight vector.

\textbf{Communication compression
methods}~\cite{alistarh2017,bernstein2018,lin2018dgc} reduce bandwidth
cost but are designed as modifications to the gradient computation step,
not as integrated components of a distributed training system with
ensemble properties and federated aggregation.

\textbf{Distributed training frameworks} such as Parameter
Server~\cite{li2014ps}, Horovod~\cite{sergeev2018}, and PyTorch
DDP~\cite{li2020ddp} provide efficient data parallelism but offer no
mechanism for ensemble diversity, non-IID data handling, or
communication compression beyond native FP16 all-reduce.

\subsection{Our Proposal}

We present \hqde{} (Hierarchical Quantized Distributed Ensemble
Learning), a framework that aims to address all three requirements jointly.
The implementation uses Ray actors when Ray is available and falls back
to local workers otherwise. Each worker keeps its model, optimizer, and
scheduler state in memory across calls. In independent mode, workers use
fixed learning-rate multipliers and dropout rates to encourage diversity.
In federated mode, workers share a reference model and synchronize at
epoch boundaries. Batches can be replicated or split across workers.
Predictions are fused with efficiency-weighted softmax or a simple mean.
In federated mode, weight deltas can be compressed with adaptive
blockwise quantization and error feedback.

\subsection{Contributions}

The principal contributions of this work are:
\begin{enumerate}
  \item To the best of our knowledge, a compact open-source implementation that combines independent ensemble training and optional FedAvg-style synchronization in a single PyTorch/Ray pipeline.
  \item An efficiency score based on per-batch accuracy and loss, used for prediction fusion and optional training aggregation.
  \item An adaptive blockwise delta quantizer with scheduled bitwidths, parameter filtering, and error feedback for communication reduction in federated mode.
  \item Configurable server aggregation (weighted mean or FedAdam) and federated normalization that can preserve batch-normalization statistics locally.
\end{enumerate}

\subsection{Paper Organization}

Section~\ref{sec:ray} introduces the Ray platform and its relevance to
distributed ensemble training. Section~\ref{sec:related} reviews related
work across ensemble methods, federated learning, communication
compression, and distributed frameworks in depth.
Section~\ref{sec:arch} describes the system architecture.
Section~\ref{sec:method} presents the technical approach with all
algorithmic details. Section~\ref{sec:setup} describes the experimental
setup. Section~\ref{sec:results} reports results with objective-based
comparison tables. Section~\ref{sec:discuss} discusses findings and
limitations. Section~\ref{sec:conc} concludes.

%% ============================================================
%%  II. THE RAY DISTRIBUTED COMPUTING PLATFORM
%% ============================================================
\section{The Ray Distributed Computing Platform}
\label{sec:ray}

\hqde{} can run with or without Ray. When Ray is available, each
ensemble member is created as a Ray actor. The actor keeps the model,
optimizer, scheduler, and efficiency score in memory across calls. The
coordinator invokes actor methods such as \texttt{train\_step},
\texttt{get\_weights}, \texttt{set\_weights}, and
\texttt{predict}, then synchronizes with \texttt{ray.get()}.
In local mode, the same calls execute on in-process workers.

\subsection{Actor-backed workers}

Ray actors provide persistent state for iterative training. The
implementation uses a thin \texttt{RayEnsembleWorker} wrapper around the
local worker class, so both Ray and local execution share the same code
path. Data is passed as PyTorch tensors per batch, which keeps the
training loop simple and consistent across backends.

\subsection{Resource hints}

Workers are created with Ray's \texttt{num\_gpus} option. The
implementation divides available GPU capacity evenly across workers, so
fractional assignments are possible. This allows multiple actors to
share one device while keeping per-worker state isolated.

\section{Related Work}
\label{sec:related}

\subsection{Ensemble Learning}

\subsubsection{Deep Ensembles}

The most widely used ensemble baseline for deep learning is Deep
Ensembles, proposed by Lakshminarayanan \etal~\cite{lakshminarayanan2017}.
The method trains $M$ neural networks from different random
initializations, each with an identical architecture and training
procedure, and combines their predictions by averaging softmax outputs:
\begin{equation}
  \bar{p}(y|\mathbf{x}) = \frac{1}{M}\sum_{m=1}^{M} p_m(y|\mathbf{x}).
  \label{eq:de}
\end{equation}
The ensemble prediction is more calibrated than any individual member
because different members make different errors. Lakshminarayanan
\etal showed that Deep Ensembles outperform Bayesian neural
networks on uncertainty quantification benchmarks, which was a
significant result given the theoretical foundations of Bayesian
methods. However, the method scales linearly in compute cost with the
number of members and assumes full dataset access at each member.

Diversity in Deep Ensembles arises entirely from random initialization
and stochastic gradient noise---the only difference between members is
their starting point and the ordering of mini-batches. This produces
relatively weak diversity. Two networks initialized with the same
architecture but different seeds may still converge to similar
regions of the loss landscape, particularly in overparameterized
settings where many global minima exist. \hqde{} addresses this by
introducing controlled hyperparameter diversity across workers,
specifically varying learning rate multipliers and dropout rates,
which drives members toward qualitatively different local minima
rather than nearby solutions in the same loss basin.

\subsubsection{Snapshot Ensembles}

Huang \etal~\cite{huang2017snapshot} proposed obtaining multiple
ensemble members from a single training run to avoid the $M\times$
compute cost of independent training. The learning rate follows a
cyclic cosine schedule that is annealed to a near-zero value and reset
multiple times; a model checkpoint is saved at each local minimum. The
final prediction averages the saved checkpoints. The schedule is:
\begin{equation}
  \alpha(t) = \frac{\alpha_0}{2}\!\left(1 + \cos\!\left(\pi
  \frac{t \bmod \lceil T/M \rceil}{\lceil T/M \rceil}\right)\right)\!,
  \label{eq:snapshot}
\end{equation}
where $\alpha_0$ is the maximum learning rate, $T$ is total iterations,
and $M$ is the number of snapshots. The limitation is that all members
share the same training trajectory; they differ only in the fine-grained
convergence behavior at each cycle's end, producing weaker diversity
than independently initialized models. In our experiments at 5 epochs
per cycle, Snapshot Ensembles achieve only $79.17\%$ on CIFAR-10 because
the architecture does not reach meaningful accuracy within a single
short cycle.

\subsubsection{TorchEnsemble}

Chen and Zhou~\cite{zhou2022torchensemble} provide a unified library of
classical ensemble methods adapted for PyTorch. Their \textbf{Voting}
classifier trains $M$ independent models on the full dataset and
combines predictions by majority vote or soft averaging. Their
\textbf{Bagging} classifier trains each member on a bootstrap sample
of the training set (approximately 63\% of the data), which reduces
per-member data but introduces dataset-level diversity. Their
\textbf{Gradient Boosting} classifier trains members sequentially, each
correcting the residuals of the previous, but this sequential structure
cannot be parallelized across workers in the way that \hqde{} requires.
TorchEnsemble methods assume full dataset access and offer no mechanism
for communication reduction or privacy preservation.

\subsubsection{MC Dropout}

Gal and Ghahramani~\cite{gal2016dropout} showed that training with
dropout and applying dropout at test time produces approximate Bayesian
inference, effectively sampling from an ensemble of models that share
all parameters except the subset activated by each dropout mask. This
is computationally efficient because it requires only one model, not
$M$ models in memory. However, the diversity generated by MC Dropout
is limited compared to truly independent networks because all members
share the same weight matrix; diversity arises only from which neurons
are masked at inference time. Accuracy improvements are typically
smaller than those from Deep Ensembles.

\subsubsection{BatchEnsemble}

Wen \etal~\cite{wen2020batchensemble} proposed BatchEnsemble, an
efficient method for creating ensemble members by multiplying the weight
matrix of each layer by rank-one perturbation matrices. Each ensemble
member $i$ has an effective weight matrix $\mathbf{W}_i =
\mathbf{r}_i\,\mathbf{s}_i^{\top} \odot \mathbf{W}$, where
$\mathbf{r}_i$ and $\mathbf{s}_i$ are learned per-member vectors.
This reduces the parameter overhead of a $M$-member ensemble from
$M\times$ to approximately $1+\epsilon$ times the base model size.
The limitation is that members still share a common weight subspace,
limiting the diversity of learned representations.

\subsubsection{Random Forests}

Breiman's Random Forests~\cite{breiman2001} established many of the
theoretical principles that motivate ensemble deep learning: training
independently on random subsets of features and data, then aggregating
predictions by majority vote. The formal treatment of bias-variance
decomposition in ensembles---that aggregating diverse, weakly correlated
predictors reduces variance without increasing bias---applies directly
to neural network ensembles. \hqde{}'s design of assigning different
hyperparameters to different workers is inspired by Random Forests'
feature randomization: the goal is not identical models with different
seeds, but genuinely different models that learn different aspects of
the data.

\subsection{Federated Learning}

\subsubsection{FedAvg}

McMahan \etal~\cite{mcmahan2017} proposed Federated Averaging (FedAvg),
the foundational algorithm for training on distributed private data.
In each round, the server broadcasts the global model to all clients.
Each client runs $E$ local epochs of SGD and returns updated weights.
The server aggregates using a sample-weighted mean:
\begin{equation}
  \mathbf{w}_{t+1} = \sum_{k=1}^{K} \frac{n_k}{n}\,\mathbf{w}_k^{(t)},
  \label{eq:fedavg}
\end{equation}
where $n_k$ is the number of samples on client $k$ and $n = \sum_k
n_k$. FedAvg's convergence analysis assumes IID data across clients;
its performance degrades under non-IID distributions because local
gradients bias toward the local optimum, causing the aggregated model
to drift from the global optimum---a phenomenon known as \textbf{client
drift}. With large $E$ and high data heterogeneity, the divergence can
be severe enough to cause the global model to perform worse than a
single client's local model.

\subsubsection{FedProx}

Li \etal~\cite{li2020fedprox} extended FedAvg to handle heterogeneous
systems and data by adding a proximal term to each client's local
objective:
\begin{equation}
  \min_{\mathbf{w}} F_k(\mathbf{w}) + \frac{\mu}{2}
  \|\mathbf{w} - \mathbf{w}^{(t)}\|^2,
  \label{eq:fedprox}
\end{equation}
where $\mu > 0$ is a proximal coefficient and $\mathbf{w}^{(t)}$ is
the current global model. The proximal term acts as a regularizer that
prevents the local model from moving too far from the global model
during local training, reducing client drift without requiring
additional communication. Empirically, FedProx provides modest
stability improvements over FedAvg under heterogeneous data but does
not close the gap with centralized training under severe non-IID
conditions.

\subsubsection{SCAFFOLD}

Karimireddy \etal~\cite{karimireddy2020scaffold} proposed SCAFFOLD to
directly correct for client drift using \textbf{control variates}. Each
client maintains a control variate $\mathbf{c}_k$ that estimates the
client's gradient bias relative to the global gradient. The local
update is:
\begin{equation}
  \mathbf{w} \leftarrow \mathbf{w} - \eta\bigl(\nabla F_k(\mathbf{w})
  - \mathbf{c}_k + \mathbf{c}\bigr),
  \label{eq:scaffold}
\end{equation}
where $\mathbf{c}$ is the server control variate. SCAFFOLD achieves
better convergence than FedAvg and FedProx under heterogeneous data
and is proven to converge regardless of data distribution. However,
it doubles the communication volume per round (transmitting both model
weights and control variates) and requires clients to store additional
state. More critically for our setting, SCAFFOLD's control variates are
designed for constant-step-size SGD; they interact destructively with
cosine annealing schedules. As the learning rate decays toward zero in
later rounds, the corrective delta $(-\mathbf{c}_k + \mathbf{c})$ is
no longer properly scaled relative to the gradient, causing client
models to diverge. This explains the unusually poor performance of
SCAFFOLD ($62.67\%$) in our experiments despite its strong theoretical
guarantees.

\subsubsection{FedBN}

Li \etal~\cite{li2021fedbn} identified batch normalization as a
specific source of degradation in federated settings. Batch
normalization layers maintain per-channel running statistics
$(\mu_{\text{run}}, \sigma^2_{\text{run}})$ that reflect the
empirical distribution of activations in the training data. When
clients have different feature distributions, their running statistics
diverge. Globally averaging these statistics produces a weighted mean
that accurately represents no individual client's data distribution,
causing feature misalignment at inference time. FedBN resolves this
by excluding batch normalization parameters from the global aggregation:
\begin{equation}
  \mathbf{w}_k^{(t+1)} = \begin{cases}
    \mathbf{w}_{\text{global}}^{(t)} &
      \text{if } \theta \notin \text{BN layers} \\
    \mathbf{w}_{k,\theta}^{(t)} &
      \text{if } \theta \in \text{BN layers}
  \end{cases}
  \label{eq:fedbn}
\end{equation}
Each client retains its own batch normalization statistics while sharing
all other parameters. This modification produces large accuracy gains
under heterogeneous feature distributions at no additional communication
cost. \hqde{} implements this as a configurable \texttt{local\_bn}
policy and shows its interaction with server optimizer choice in
the ablation study.

\subsubsection{FedAdam and Adaptive Server Optimizers}

Reddi \etal~\cite{reddi2021fedadam} proposed applying adaptive
server-side optimizers---FedAdam, FedYogi, and FedAdagrad---to the
aggregated client updates. The motivation is that different parameters
may converge at different rates across rounds; adaptive scaling on the
server side should accelerate slow-converging parameters while
stabilizing fast-converging ones. The FedAdam update is:
\begin{equation}
  \mathbf{w}_{t+1} = \mathbf{w}_t + \eta_s
  \frac{\hat{\mathbf{m}}_t}{\sqrt{\hat{\mathbf{v}}_t} + \varepsilon},
  \label{eq:fedadam}
\end{equation}
where $\hat{\mathbf{m}}_t$ and $\hat{\mathbf{v}}_t$ are bias-corrected
first and second moment estimates of the aggregated gradient. Reddi
\etal demonstrate convergence improvements over FedAvg in several
settings. However, their analysis assumes that moment estimates are
formed from a sufficient number of rounds to be reliable. In our
setting with only 4 workers and 20 rounds, moment estimates are noisy
and the adaptive scaling amplifies rather than dampens early instability.
Furthermore, as we show, the interaction between FedAdam and shared
batch normalization can create a severe failure mode in federated
settings.

\subsubsection{FedNTD}

Lee \etal~\cite{lee2022fedndt} proposed Federated Not-True Distillation
(FedNTD) to address global knowledge forgetting in federated learning.
When a client trains on its local data partition, it sees only a
subset of the global class distribution. For classes absent from the
local data, the local model's predictions are not directly optimized
and may drift from the global model's learned representation. FedNTD
adds a knowledge distillation term to the local objective that penalizes
divergence between the local model's predictions and the global model's
predictions on the non-local classes:
\begin{equation}
  \mathcal{L}_{\text{NTD}} = \mathcal{L}_{\text{CE}} + \lambda
  \cdot \mathrm{KL}\!\left(\bar{p}_k^{\neg y}(\mathbf{x}) \,\|\,
  \bar{p}_g^{\neg y}(\mathbf{x})\right),
  \label{eq:fedndt}
\end{equation}
where $\bar{p}_k^{\neg y}$ and $\bar{p}_g^{\neg y}$ are the local and
global predictions excluding the true class. This prevents local
models from catastrophically forgetting global knowledge at the cost
of an additional forward pass through the global model per batch.

\subsubsection{Personalized Federated Learning}

Several works address federated settings where a single global model
is insufficient because client distributions are fundamentally different.
pFedMe~\cite{dinh2020pfedme} frames personalized FL as a bilevel
optimization problem: each client optimizes a personalized model that
remains close to the global model. Ditto~\cite{li2021ditto} maintains
both a global and a local model per client, training the local model
with a regularization term toward the global model. These methods
produce per-client models rather than a single global model, which is
conceptually related to \hqde{}'s ensemble design. However, they focus
on personalization rather than ensemble accuracy improvement, and they
do not provide communication compression or the efficiency-weighted
prediction fusion mechanism.

\subsection{Communication Compression}

\subsubsection{Gradient Quantization: QSGD}

Alistarh \etal~\cite{alistarh2017} proposed Quantized SGD (QSGD) as
the first rigorous treatment of gradient quantization with convergence
guarantees. Each gradient element $g_i$ is stochastically quantized to
one of $s+1$ levels by:
\begin{equation}
  \mathcal{Q}(g_i) = \|g\|_2 \cdot \mathrm{sgn}(g_i) \cdot
  \xi\!\left(\frac{|g_i|}{\|g\|_2}, s\right),
  \label{eq:qsgd}
\end{equation}
where $\xi(\cdot, s)$ is a stochastic rounding function and $s$ controls
the number of quantization levels. QSGD is an unbiased estimator of
the true gradient: $\mathbb{E}[\mathcal{Q}(g)] = g$. The variance of
the quantized gradient is bounded, enabling proof of convergence at
rates similar to unquantized SGD. QSGD achieves $2\times$ to $8\times$
compression depending on the choice of $s$. A key difference from
\hqde{}'s quantization is that QSGD operates on gradients (before
weight aggregation) and requires modification of the backward pass,
while \hqde{} operates on weight deltas (after local training) and
is transparent to the client optimizer.

\subsubsection{Sign-Based Compression: SignSGD}

Bernstein \etal~\cite{bernstein2018} proposed transmitting only the
sign of each gradient element, achieving $32\times$ compression by
reducing each float32 to a single bit. The server aggregates using
majority vote: $w_{t+1} = w_t - \eta\,\mathrm{sign}(\sum_k
\mathrm{sign}(\nabla F_k))$. SignSGD achieves linear speedup with
the number of workers under certain conditions but loses all magnitude
information, which limits accuracy on problems requiring precise
gradient directions. Under non-IID data distributions, majority vote
does not provide unbiased gradient estimates, causing convergence
degradation.

\subsubsection{Sparsification: Top-$k$ and Deep Gradient Compression}

Gradient sparsification methods transmit only the $k$ gradient
components with the largest magnitude and set the rest to zero.
Aji and Heafield~\cite{aji2017} showed that transmitting only 1\% of
gradient components can reduce communication volume by $99\times$ with
minimal accuracy loss. Lin \etal~\cite{lin2018dgc} proposed Deep
Gradient Compression (DGC), which combines Top-$k$ sparsification with
momentum correction, local gradient accumulation, and warmup training
to prevent accuracy degradation. The momentum correction term adds
locally accumulated gradients that were not transmitted in previous
rounds to the current gradient before sparsification. This is
functionally similar to the error feedback mechanism in \hqde{}'s
quantization, but applied to gradient sparsification rather than
weight delta quantization.

\subsubsection{Low-Rank Approximation: PowerSGD}

Vogels \etal~\cite{vogels2019powersgd} proposed PowerSGD, which
approximates gradient matrices as low-rank products $\mathbf{G}
\approx \mathbf{P}\mathbf{Q}^{\top}$ where $\mathbf{P}$ and
$\mathbf{Q}$ have rank $r \ll \min(m,n)$. The low-rank factors are
transmitted instead of the full gradient matrix, achieving compression
ratios of $r\times$ to $20\times$ depending on the gradient structure.
PowerSGD performs well on layers with approximately low-rank gradients
(such as linear layers in transformers) but provides less benefit on
convolutional layers whose gradients may not have a naturally low-rank
structure.

\subsubsection{Error Feedback Theory}

Karimireddy \etal~\cite{karimireddy2019ef} established the theoretical
foundation for error feedback (EF) in compressed communication.
Without error feedback, gradient compression introduces a bias that
can prevent convergence or cause the model to converge to a
non-optimal point. With error feedback, the compression error from
each round is accumulated and added to the gradient in the next round,
ensuring that all gradient information is eventually transmitted. EF
is a necessary component of any compression scheme applied to
non-unbiased compressors (including Top-$k$ and deterministic
quantization). \hqde{}'s quantization module implements EF via a
per-worker residual buffer $\mathbf{r}_k^{(t)}$, ensuring that
quantization-induced bias does not accumulate across rounds.

\subsubsection{Adaptive Precision: 1-bit Adam}

Tang \etal~\cite{tang2021onebitadam} proposed 1-bit Adam, which
combines Adam's adaptive moment correction on the client with 1-bit
gradient compression during the communication step. The insight is
that once Adam's variance estimate stabilizes after a warmup period,
the effective gradient direction changes slowly and can be compressed
aggressively without information loss. This motivates the scheduled
bitwidth decay in \hqde{}: early rounds use more bits (higher precision)
because weight deltas are large and informative; later rounds can use
fewer bits because deltas have stabilized and the remaining gradient
information is less critical for accuracy.

\subsection{Distributed Training Infrastructure}

\subsubsection{Parameter Server Architecture}

Li \etal~\cite{li2014ps} proposed the parameter server (PS)
architecture, in which one or more servers maintain the global model
and multiple worker nodes compute gradients. Workers push gradient
updates to servers and pull updated parameters. The PS architecture
supports asynchronous updates (where workers do not wait for each other)
and synchronous updates (where the server waits for all workers before
aggregating). Asynchronous PS can suffer from \textbf{staleness}: a
worker may compute a gradient based on a model that has already been
updated multiple times by other workers, causing the applied gradient
to be misaligned with the current model. \hqde{} uses synchronous
aggregation only, eliminating staleness at the cost of requiring all
workers to complete each round before proceeding.

\subsubsection{Ring-AllReduce and Horovod}

Baidu proposed ring-allreduce~\cite{baidu2017ring} as an alternative
to the parameter server architecture for synchronous data-parallel
training. Workers are arranged in a virtual ring; each worker
simultaneously sends its gradient shard to the next worker and receives
a gradient shard from the previous worker. After $n-1$ steps, each
worker holds the sum of all workers' gradients for a particular shard.
This is bandwidth-optimal: total communication volume equals $2 \times
(n-1)/n \times \text{model size}$, independent of the number of
workers. Horovod~\cite{sergeev2018} made ring-allreduce accessible for
TensorFlow, PyTorch, and MXNet through a simple API. Horovod is
efficient for homogeneous cluster training but requires all workers to
participate in each communication round and does not support the
ensemble diversity or federated aggregation modes needed by \hqde{}.

\subsubsection{PyTorch Distributed Data Parallel}

Li \etal~\cite{li2020ddp} describe PyTorch's Distributed Data Parallel
(DDP) module, which replicates the model on each device and
synchronizes gradients using NCCL's efficient all-reduce implementation
after each backward pass. DDP uses \textbf{gradient bucketing}: rather
than transmitting one gradient at a time, gradients are accumulated
into buckets that are all-reduced together when full, amortizing
communication overhead over multiple parameters. DDP achieves near-linear
scaling efficiency on homogeneous clusters with fast interconnects
(NVLink or InfiniBand). Its limitations are: (1) all workers must have
identical batch sizes, (2) data must be IID across workers, (3) all
workers train the same model, preventing ensemble diversity, and
(4) communication compression beyond native FP16 requires third-party
libraries that do not integrate cleanly with the existing API.

\subsection{Quantization Methods}

\subsubsection{Post-Training Quantization}

Post-training quantization (PTQ) reduces the precision of model
weights and activations after training is complete, without any
retraining. Symmetric uniform quantization maps floating-point values
to integers using a scale factor derived from the maximum absolute
value of the tensor. PTQ can reduce model size by $4\times$ (float32
to int8) with minimal accuracy loss on many classification tasks.
However, for very low bitwidths (below 4 bits), PTQ accuracy degrades
significantly. \hqde{}'s blockwise quantization applies PTQ-style
symmetric quantization to weight deltas rather than absolute weights,
which is more effective because delta magnitudes are smaller and more
uniformly distributed than raw weight magnitudes.

\subsubsection{Quantization-Aware Training}

Quantization-aware training (QAT) simulates quantization during the
forward pass using \textbf{fake quantization} operations: weights are
quantized and immediately dequantized before use, so the model learns
to be robust to quantization error. QAT produces more accurate
quantized models than PTQ at low bitwidths because the model adjusts
its weight distribution during training to minimize the impact of
quantization. \hqde{} does not apply QAT to the model itself but uses
a related principle for communication: the quantization error
(residual) is accumulated and added to the next round's delta,
which functions as implicit QAT at the communication level.

\subsubsection{Mixed-Precision Training}

Micikevicius \etal~\cite{micikevicius2018mixed} proposed training with
FP16 activations and FP32 master weights, using loss scaling to prevent
gradient underflow. The forward and backward passes run in FP16,
reducing memory consumption and increasing throughput on hardware with
tensor cores; the weight update is applied to a separate FP32 copy.
\hqde{} uses mixed-precision training (via PyTorch's
\texttt{torch.cuda.amp}) for the local training step but transmits
weight deltas in quantized integer format for communication efficiency.
These two precision-reduction mechanisms operate at different levels:
AMP reduces training-time memory and compute, while delta quantization
reduces communication volume.

%% ============================================================
%%  IV. SYSTEM ARCHITECTURE
%% ============================================================
\section{System Architecture}
\label{sec:arch}

\subsection{Architectural Overview}

\hqde{} is organized into four functional subsystems as shown in
Fig.~\ref{fig:arch}. The \textbf{Coordinator} (the \texttt{HQDESystem}
and \texttt{DistributedEnsembleManager} classes) manages the training
lifecycle: it creates workers, runs the training loop, and broadcasts
updated weights. The \textbf{Worker Ensemble} consists of $K$ Ray actor
processes when Ray is available, or local worker instances otherwise.
Each worker holds its own model, optimizer, scheduler, and efficiency
score in memory. The \textbf{Aggregation and Quantization Module} is
enabled only in federated mode; it computes weight deltas, optionally
compresses them with adaptive blockwise quantization, and aggregates them
with a weighted mean or a FedAdam-style server update. The
\textbf{Prediction Fusion Module} combines worker logits using
efficiency-weighted softmax or a simple mean. Optional modules for
hierarchical aggregation and Byzantine fault tolerance exist in the
repository but are not wired into the core training loop.

\begin{figure}[!t]
\centering
\vspace{4.5cm}
\caption{High-level architecture of \hqde{}. The Coordinator initializes
$K$ workers (Ray actors when available) and drives the batch-level
training loop. In federated mode, workers transmit weight deltas after
each epoch, optional quantization is applied, and the server aggregates
and broadcasts updated reference weights. At inference, the Prediction
Fusion Module combines worker logits using efficiency-weighted or mean
aggregation.}
\label{fig:arch}
\end{figure}

\subsection{Backbone Architecture: Modified ResNet-18}

Standard ResNet-18~\cite{he2016resnet} begins with a $7\times7$
convolution at stride 2 followed by a $3\times3$ max-pooling layer at
stride 2. Together these produce a spatial downsampling factor of 4
before the first residual block. For a $224\times224$ input this
reduces the feature map to $56\times56$, preserving adequate spatial
resolution. For $32\times32$ inputs---used in CIFAR-10 and related
benchmarks---the same sequence reduces the feature map to $8\times8$
after just two layers, before any residual computation has occurred.
Spatial information critical for distinguishing objects at this
resolution is discarded in the first few milliseconds of the forward
pass.

\hqde{} uses a modified backbone, as illustrated in Fig.~\ref{fig:resnet},
with three targeted changes. First, the $7\times7$ convolution at
stride 2 is replaced with a $3\times3$ convolution at stride 1, which
extracts features without downsampling. Second, the initial max-pooling
layer is removed entirely. Third, a configurable dropout layer is
inserted before the final linear classifier, enabling per-worker
regularization diversity at different dropout rates. All four residual
block groups remain unchanged. The total parameter count is preserved
at $11.17\times10^6$.

\begin{figure}[!t]
\centering
\vspace{4cm}
\caption{Comparison of (a)~standard ResNet-18 entry sequence for
$224\times224$ inputs and (b)~the \hqde{} modified backbone for
$32\times32$ inputs. The $7\times7$ stride-2 convolution is replaced
by a $3\times3$ stride-1 convolution; the initial max-pool is removed;
a per-worker dropout is added before the classifier. Parameter count
($11.17\times10^6$) is unchanged.}
\label{fig:resnet}
\end{figure}

\subsection{Training Modes}

\hqde{} supports two primary training modes.

\textbf{Independent Ensemble Mode.} Each of the $K$ workers trains on
the full dataset with distinct hyperparameters: learning rate
multipliers $\lambda_k \in \{1.0,\,0.85,\,1.15,\,0.95\}$ and dropout
rates $d_k \in \{0.10,\,0.12,\,0.15,\,0.13\}$. Workers never
synchronize weights during training, so diversity arises from different
optimization trajectories.

\textbf{Federated Mode.} Workers start from a shared reference model and
synchronize weight deltas after each epoch. The library expects the
caller to provide the data partitioning; within each step, batches are
either replicated or split across workers according to the
\texttt{batch\_assignment} setting. Hyperparameter diversity is disabled
in this mode to keep update scales consistent during aggregation.

%% ============================================================
%%  V. TECHNICAL APPROACH
%% ============================================================
\section{Technical Approach}
\label{sec:method}

\subsection{Efficiency-Weighted Ensemble Prediction}

After each training batch, worker $k$ computes a performance signal:
\begin{equation}
  S_t^{(k)} = 2\,\mathrm{acc}_t^{(k)} + \frac{1}{1+\mathrm{loss}_t^{(k)}}.
  \label{eq:signal}
\end{equation}
The signal weights accuracy at twice the importance of loss; the
reciprocal loss term is bounded between 0 and 1, preventing
loss-dominated early training from suppressing accurate workers.
The running efficiency score is:
\begin{equation}
  e_t^{(k)} = 0.8\,e_{t-1}^{(k)} + 0.2\,S_t^{(k)}.
  \label{eq:ema}
\end{equation}
The exponential moving average with coefficient 0.8 smooths batch
variance while remaining responsive to sustained accuracy changes. At
inference time:
\begin{equation}
  f(\mathbf{x}) = \sum_{k=1}^K w_k\,f_k(\mathbf{x}), \quad
  w_k = \frac{\exp(e^{(k)})}{\sum_{j=1}^K \exp(e^{(j)})}.
  \label{eq:pred_agg}
\end{equation}
Softmax normalization ensures the weights sum to one and amplifies
differences between workers' efficiency scores. The same scores can
optionally weight delta aggregation during training.

\subsection{Adaptive Blockwise Delta Quantization}

\hqde{} transmits weight deltas rather than absolute weights when
federated mode is enabled. The delta for worker $k$ at round $t$ is:
\begin{equation}
  \Delta\mathbf{w}_k^{(t)} = \mathbf{w}_k^{(t)} - \mathbf{w}_{\mathrm{ref}}.
  \label{eq:delta}
\end{equation}
The flattened delta tensor is split into contiguous blocks of
$B = 1\,024$ elements. The code computes an elementwise importance score
from normalized absolute values:
\begin{equation}
  I = \frac{\lvert\Delta\mathbf{w}\rvert - \min\lvert\Delta\mathbf{w}\rvert}
           {\max\lvert\Delta\mathbf{w}\rvert - \min\lvert\Delta\mathbf{w}\rvert}.
  \label{eq:importance}
\end{equation}
When the denominator is zero, $I$ is set to $0.5$. Each block uses the
mean importance $\bar{I}_i$ to choose a bitwidth. The round-level target
bitwidth $b_t$ is scheduled by phase: 32-bit during warmup rounds, then
$b_{\max}$, $b_{\mathrm{base}}$, and $b_{\min}$ over the early, middle,
and late thirds of the remaining rounds. For block $i$:
\begin{equation}
  b_i = \mathrm{clamp}\Bigl(\mathrm{round}((b_t - 2) + 4\,\bar{I}_i),
  b_{\min}, b_{\max}\Bigr).
  \label{eq:bits}
\end{equation}
Symmetric uniform quantization then uses:
\begin{align}
  s_i &= \frac{\max\lvert\Delta\mathbf{w}_i\rvert}{2^{b_i-1}-1},
  \label{eq:scale}\\[3pt]
  \Delta\hat{\mathbf{w}}_i &= \mathrm{clamp}\Bigl(
    \mathrm{round}(\Delta\mathbf{w}_i / s_i),
    -q_{\max}, q_{\max}\Bigr),
  \label{eq:quant}
\end{align}
where $q_{\max} = 2^{b_i-1}-1$ and the dequantized block is
$\Delta\hat{\mathbf{w}}_i \cdot s_i$. If error feedback is enabled,
the residual is stored per worker and added to the next round's delta.
The quantizer skips non-floating tensors, small tensors, bias terms, and
normalization parameters.

\subsection{Server-Side Aggregation}

The coordinator aggregates deltas, not full weights. The default update is
a weighted mean:
\begin{equation}
  \Delta\bar{\mathbf{w}} = \sum_{k=1}^K \alpha_k \Delta\mathbf{w}_k,
  \quad \sum_{k=1}^K \alpha_k = 1,
  \label{eq:agg}
\end{equation}
\begin{equation}
  \mathbf{w}_{t+1} = \mathbf{w}_{\mathrm{ref}} + \Delta\bar{\mathbf{w}}.
  \label{eq:agg_update}
\end{equation}
The weights $\alpha_k$ come from per-epoch sample counts
(\texttt{sample\_weighted}), a uniform mean, or efficiency scores
(\texttt{efficiency\_weighted}).

\textbf{FedAdam}~\cite{reddi2021fedadam} treats $\Delta\bar{\mathbf{w}}$
as the server update and maintains first and second moments. The code
applies bias correction and then updates the reference weights using the
server learning rate in Equation~(\ref{eq:fedadam}).

\subsection{Batch Normalization Policy}

In \textbf{local BN} mode, batch normalization parameters (running
mean, running variance, scale, shift) are excluded from the global
broadcast. Each worker retains locally-calibrated statistics. In
\textbf{shared BN} mode, all parameters including BN are synchronized.
This option is only applied in federated mode.
The critical interaction between this choice and the server optimizer is
analyzed in Section~\ref{sec:ablation}.

\subsection{Hierarchical Aggregation and Fault Tolerance}

The repository includes optional aggregation utilities that are not wired
into the \texttt{HQDESystem} training loop. The
\texttt{HierarchicalAggregator} builds a tree of Ray actors and performs
weighted-mean aggregation bottom-up. The
\texttt{ByzantineFaultTolerantAggregator} provides outlier filtering
(MAD, cosine similarity, or Euclidean distance) and falls back to a
geometric median when reliable sources are scarce. These modules can be
invoked independently when hierarchical or robust aggregation is
required.

\subsection{Training Algorithm}

Algorithm~\ref{alg:hqde} presents the complete \hqde{} training
procedure.

\begin{algorithm}[!t]
\small
\SetAlgoLined
\DontPrintSemicolon
\KwIn{Dataloader $\mathcal{D}$, workers $K$, epochs $T$, ensemble\_mode,
      batch\_assignment, server\_opt, bn\_mode, quant\_config}
\KwOut{Ensemble models $\{M_k\}_{k=1}^K$}

\BlankLine
\tcp{--- Initialization ---}
Create $K$ workers and initialize optimizers and schedulers\;
Set efficiency scores $e_k \leftarrow 1$\;
\If{ensemble\_mode $=$ fedavg}{
  Broadcast initial weights from worker 0 to all workers\;
  Set reference weights $\mathbf{w}_{\mathrm{ref}}$\;
  Reset quantizer state\;
}

\BlankLine
\tcp{--- Training epochs ---}
\For{$t = 1$ \KwTo $T$}{
  Reset per-epoch sample counts $n_k \leftarrow 0$\;
  \For{each batch $(\mathbf{X}, \mathbf{y})$ in $\mathcal{D}$}{
    Build worker batches by \texttt{batch\_assignment}\;
    \For{$k = 1$ \KwTo $K$ \textup{(parallel)}}{
      Run \texttt{train\_step} on worker $k$\;
      Update $e_k$ inside the worker\;
      Accumulate $n_k$\;
    }
  }
  \If{ensemble\_mode $=$ fedavg}{
    Collect weights and compute $\Delta\mathbf{w}_k = \mathbf{w}_k - \mathbf{w}_{\mathrm{ref}}$\;
    \If{quant\_config enabled and $t$ $>$ warmup}{
      Quantize $\Delta\mathbf{w}_k$ with error feedback\;
    }
    Aggregate deltas with sample or efficiency weights\;
    Apply \texttt{server\_opt} (mean or FedAdam)\;
    Broadcast updated weights (skip BN keys if \texttt{bn\_mode} = local)\;
    Update $\mathbf{w}_{\mathrm{ref}}$\;
  }
  Step per-worker schedulers\;
}

\BlankLine
\tcp{--- Inference ---}
\Return $f(\mathbf{x}) = \sum_k w_k f_k(\mathbf{x})$ with
$w_k \propto \exp(e_k)$ or mean weights\;
\caption{\hqde{} Training Procedure}
\label{alg:hqde}
\end{algorithm}

\subsection{Implementation Details}

First, the coordinator calls \texttt{create\_ensemble\_workers}. When Ray
is available, each worker is a \texttt{RayEnsembleWorker} actor created
with \texttt{num\_gpus} set to the available GPU count divided by the
number of workers. Otherwise a local worker instance is created. Each
worker selects a CUDA device when available, uses \texttt{channels\_last}
format for image batches, and can enable \texttt{torch.compile}.

Next, training uses \texttt{train\_ensemble}. For each batch, the
coordinator builds per-worker batches by replication or splitting. It
invokes \texttt{train\_step} on each worker and collects loss, accuracy,
and sample counts. The worker updates its efficiency score after every
batch using Equation~(\ref{eq:ema}).

Then, in federated mode, the coordinator collects state dictionaries and
computes deltas against the reference weights. The
\texttt{AdaptiveQuantizer} skips non-floating tensors, small tensors,
bias terms, and normalization parameters. It performs blockwise
symmetric quantization, updates residual buffers when error feedback is
enabled, and records compression metrics.

Then, aggregated deltas are combined with sample-weighted or
efficiency-weighted averaging. The server update uses a weighted mean or
FedAdam with stored first and second moments. The coordinator broadcasts
the updated weights and optionally preserves local batch-normalization
keys.

Finally, prediction uses \texttt{predict\_batch} with mean or
efficiency-weighted softmax. Training uses AMP, gradient clipping, and a
warmup + cosine learning-rate schedule in each worker.

%% ============================================================
%%  VI. EXPERIMENTAL SETUP
%% ============================================================
\section{Experimental Setup}
\label{sec:setup}

\subsection{Datasets}

\textbf{CIFAR-10}~\cite{krizhevsky2009} is the primary dataset used in
this repository's example scripts. The examples use the transforms in
\texttt{DataPreprocessor.cifar10\_transforms}: random crop (padding~4)
and random horizontal flip at training time, followed by normalization.
The core training loop accepts any PyTorch dataset and dataloader, so
additional datasets can be evaluated with the same pipeline, but those
experiments are not bundled in the codebase.

\subsection{Federated Partitioning}

The library does not implement dataset partitioning. Users provide their
own non-IID splits or rely on \texttt{batch\_assignment=split} to divide
mini-batches across workers. The training loop records per-epoch sample
counts and uses them as aggregation weights when
\texttt{training\_aggregation} is set to \texttt{sample\_weighted}.

\subsection{Training Hyperparameters}

The CIFAR-10 example script uses \texttt{make\_cifar\_training\_config}:
SGD with Nesterov momentum $0.9$, weight decay $5\times10^{-4}$, initial
learning rate $0.1$, cosine annealing with 5-epoch linear warmup
($0.2\times \to 1.0\times$), label smoothing $0.1$, gradient clip at norm
$1.0$, AMP (FP16), and seed 42. The script uses batch size 128 for
training and 256 for evaluation. The number of epochs or rounds is
configurable in the training loop.

FedAdam parameters: server LR $1.0$, $\beta_1=0.9$, $\beta_2=0.99$,
$\varepsilon=10^{-3}$.

Default quantization: $b_{\mathrm{base}}=8$, $b_{\min}=4$,
$b_{\max}=16$, block size $B=1\,024$, 5-round warmup.

\subsection{Hardware}

Single-GPU: NVIDIA L4 (22~GB VRAM) on Google Colab.
Multi-GPU: $2\times$ NVIDIA Tesla T4 (14.6~GB each) on Kaggle.

\subsection{Baselines}

\textbf{Distributed:} Centralized single-model; PyTorch DDP (1 GPU).

\textbf{Ensemble:} Deep Ensembles ($K=4$); TorchEnsemble Voting and
Bagging ($K=4$, 5 epochs); Snapshot Ensembles (5 epochs/cycle).

\textbf{Federated:} FedAvg~\cite{mcmahan2017};
FedProx~\cite{li2020fedprox} ($\mu=0.01$);
SCAFFOLD~\cite{karimireddy2020scaffold};
FedNTD~\cite{lee2022fedndt}.

All baselines use the same modified ResNet-18 and identical augmentation.

%% ============================================================
%%  VII. RESULTS AND ANALYSIS
%% ============================================================
\section{Results and Analysis}
\label{sec:results}

The results below summarize example runs of the \hqde{} implementation.
Baseline numbers are reported for context and were produced with
separate reference implementations. Results are organized into five
objective-based comparisons: (A) vs.\ distributed and ensemble
baselines; (B) vs.\ federated baselines; (C) communication efficiency;
(D) ablation; (E) multi-GPU runs; and (F) multi-dataset generalization.

\subsection{Accuracy vs.\ Distributed and Ensemble Baselines}
\label{sec:vs_ensemble}

Table~\ref{tab:vs_ensemble} shows that \hqde{} Independent mode achieves
$94.27\%$ test accuracy, surpassing centralized training by $1.15$~pp
and DDP by the same margin. The gain over Deep Ensembles is $6.86$~pp
($94.27\%$ vs.\ $87.41\%$) despite identical model architectures and
equivalent total compute. This gap is substantially larger than typical
random-seed ensemble gains ($\approx1$--$2$~pp) and reflects the
contribution of deliberate hyperparameter diversity. Deep Ensembles
assigns identical hyperparameters to all members; diversity arises only
from random initialization noise and stochastic gradient ordering.
\hqde{}'s workers are assigned different learning rate multipliers and
dropout rates, driving them toward qualitatively different regions of
the loss landscape and producing stronger complementary representations.

Snapshot Ensembles achieve $79.17\%$ because 5 epochs per cosine cycle
is insufficient for the modified ResNet-18 to reach its accuracy peak
at each local minimum. All members share the same optimization
trajectory up to the checkpoint; diversity is minimal. TorchEnsemble
Bagging ($80.07\%$) trains each member on $\approx63\%$ of the data
per bootstrap sample, limiting accuracy relative to full-data training.

\textbf{Memory efficiency.} \hqde{} workers operate at $17.2$~MB peak
GPU memory per process because Ray actors own their model in process-local
memory and the main process holds no model copy. Deep Ensembles requires
$735.4$~MB ($43\times$ more) because all four models coexist in GPU
memory during training. TorchEnsemble Bagging requires $1\,115.2$~MB
($65\times$ more).

\begin{table}[!t]
\caption{Comparison Against Distributed and Ensemble Baselines.
(CIFAR-10, NVIDIA L4, 20 epochs, $K=4$).
$\Delta$ = difference from centralized baseline.}
\label{tab:vs_ensemble}
\centering
\small
\renewcommand{\arraystretch}{1.25}
\setlength{\tabcolsep}{3.5pt}
\begin{tabular}{L{2.6cm}C{1.2cm}C{1.2cm}C{1.2cm}C{1.5cm}}
\toprule
\textbf{Method} & \textbf{Acc.} & \textbf{$\Delta$} & \textbf{Time} & \textbf{Mem.} \\
 & (\%) & (pp) & (s) & (MB) \\
\midrule
\multicolumn{5}{l}{\textit{Distributed}} \\
Centralized            & 93.12 & ---     & \textbf{204} & 473.0 \\
DDP (1 GPU)            & 93.12 & $0.00$  & \textbf{204} & 514.3 \\
\midrule
\multicolumn{5}{l}{\textit{Ensemble}} \\
Deep Ensemble          & 87.41 & $-5.71$ & \textbf{204} & 735.4 \\
TE Bagging (5 ep.)     & 80.07 & $-13.05$& 376          & 1\,115.2 \\
TE Voting (5 ep.)      & 79.37 & $-13.75$& 376          & 1\,119.7 \\
Snapshot (5 ep./cyc.)  & 79.17 & $-13.95$& 54           & 646.0 \\
\midrule
\multicolumn{5}{l}{\textit{\hqde{} (this work)}} \\
\hqde{} Independent
  & \textbf{94.27} & \textbf{$+1.15$} & 643 & \textbf{17.2} \\
\bottomrule
\end{tabular}
\end{table}

\subsection{Accuracy vs.\ Federated Baselines}
\label{sec:vs_federated}

Table~\ref{tab:vs_federated} compares \hqde{} federated configurations
against all federated baselines on a non-IID CIFAR-10 split prepared
outside the library.
\hqde{} FedAvg with the weighted mean server achieves $89.41\%$, and
adding 4-to-16 bit quantization raises this to $89.62\%$. Both results
exceed all standard federated algorithms: FedAvg ($86.77\%$), FedProx
($85.87\%$), FedNTD ($85.60\%$), and SCAFFOLD ($62.67\%$).

FedAvg outperforms FedProx ($86.77\%$ vs.\ $85.87\%$) in this setting.
The proximal term in FedProx constrains local updates, which is
beneficial when client drift is severe. Under cosine annealing, however,
the learning rate decays naturally in later rounds, reducing drift
without the proximal term. Adding the proximal constraint on top of a
decaying learning rate over-regularizes local training.

FedNTD ($85.60\%$) performs slightly below FedProx. The KL divergence
distillation term requires an additional forward pass through the global
model per batch, increasing per-round compute. With only 1 local epoch
per round, this overhead does not translate to additional gradient
updates, effectively reducing training efficiency relative to simpler
methods.

\textbf{SCAFFOLD failure mode.} SCAFFOLD achieves $62.67\%$, the
lowest result in the entire comparison. As explained in
Section~\ref{sec:related}, SCAFFOLD's control variates are calibrated
under the assumption of a constant step size. The cosine schedule
drives the learning rate toward zero in later rounds; the control
variate corrections, unchanged in scale, dominate the gradient and
cause client models to diverge. A partial recovery occurs in the final
few rounds as the corrections begin to align with the diminishing
gradient, but insufficient accuracy is reclaimed.

\begin{table}[!t]
\caption{Comparison Against Federated Learning Baselines.
(CIFAR-10, non-IID split, NVIDIA L4, 20 rounds, $K=4$).
$\Delta$ = difference from FedAvg baseline.}
\label{tab:vs_federated}
\centering
\small
\renewcommand{\arraystretch}{1.25}
\setlength{\tabcolsep}{3.5pt}
\begin{tabular}{L{2.6cm}C{1.2cm}C{1.2cm}C{1.2cm}C{1.5cm}}
\toprule
\textbf{Method} & \textbf{Acc.} & \textbf{$\Delta$} & \textbf{Time} & \textbf{Mem.} \\
 & (\%) & (pp) & (s) & (MB) \\
\midrule
\multicolumn{5}{l}{\textit{Federated baselines}} \\
FedAvg         & 86.77 & ---       & 352 & 758.4 \\
FedProx        & 85.87 & $-0.90$   & 372 & 870.8 \\
FedNTD         & 85.60 & $-1.17$   & 447 & 934.1 \\
SCAFFOLD       & 62.67 & $-24.10$  & 575 & 802.1 \\
\midrule
\multicolumn{5}{l}{\textit{\hqde{} federated (this work)}} \\
\hqde{} (FedAdam)  & 80.82 & $-5.95$  & 431 & \textbf{17.2} \\
\hqde{} (Mean)     & 89.41 & $+2.64$  & 428 & \textbf{17.2} \\
\hqde{} + Quant    & \textbf{89.62} & \textbf{$+2.85$} & 429 & \textbf{17.2} \\
\bottomrule
\end{tabular}
\end{table}

\subsection{Communication Efficiency}
\label{sec:comm}

Table~\ref{tab:quant} reports communication volume and accuracy across
quantization configurations. The standard 4-to-16 bit configuration
compresses total communication from $4\,473.4$~MB to $565.9$~MB
($7.9\times$) with a marginal accuracy \emph{gain} of $+0.21\%$. The
accuracy improvement is reproducible across runs. Its mechanism is
quantization noise acting as an implicit regularizer on weight delta
updates, preventing any single round's large gradient components from
dominating the global model update.

Aggressive 2-to-8 bit quantization achieves $12.6\times$ compression
at $80.04\%$ accuracy---a $0.78\%$ reduction from the uncompressed
baseline. Despite using only 2--8 bits per block, the error feedback
mechanism recovers most of the gradient information across subsequent
rounds, keeping accuracy degradation below $1\%$ even at this extreme
compression ratio.

\begin{table}[!t]
\caption{Communication Efficiency: Quantization Comparison.
(\hqde{} FedAvg, Mean Server, CIFAR-10, 20 rounds.)}
\label{tab:quant}
\centering
\small
\renewcommand{\arraystretch}{1.25}
\setlength{\tabcolsep}{4pt}
\begin{tabular}{lcccc}
\toprule
\textbf{Configuration} & \textbf{Bits} & \textbf{Acc.} (\%) & \textbf{Comm.} (MB) & \textbf{Ratio} \\
\midrule
No quantization  & 32   & 89.41          & 4\,473.4 & 1.0$\times$ \\
Standard         & 4--16 & \textbf{89.62} & 565.9    & \textbf{7.9$\times$} \\
Aggressive       & 2--8  & 80.04          & 354.7    & 12.6$\times$ \\
\bottomrule
\end{tabular}
\end{table}

\subsection{Ablation Study}
\label{sec:ablation}

Table~\ref{tab:ablation} presents the full ablation. The most critical
finding is the interaction between the server optimizer and the batch
normalization policy.

\textbf{Server optimizer has the largest effect on accuracy.} Switching from FedAdam
to the weighted mean improves accuracy by $8.59$~pp ($80.82\%
\to 89.41\%$). At $K=4$ and $T=20$, FedAdam's second moment estimate
$\hat{v}_t$ is formed from insufficient observations to be reliable.
Early-round moment estimates are dominated by noise; the adaptive
scaling $1/\sqrt{\hat{v}_t}$ amplifies unstable early gradient
components rather than normalizing them. The client-side cosine
schedule with warmup already provides well-scaled updates, making
server-side adaptive scaling redundant at this scale.

\textbf{BN drop under FedAdam.} The combination of FedAdam and
shared BN reduces accuracy to $10.00\%$---random chance for 10 classes.
The mechanism is gradient-scale mismatch: Adam's update magnitude is
proportional to $g/\sqrt{\hat{v}}$ where $\hat{v}$ is the moving
average of squared gradients. Batch normalization parameters (scale,
shift, running mean, running variance) have gradient magnitudes
orders-of-magnitude smaller than convolutional weight gradients.
When BN parameters are included in the globally broadcast payload,
the server's $\hat{v}$ is dominated by convolutional gradient magnitudes.
The resulting $1/\sqrt{\hat{v}}$ for BN parameters is enormous,
producing updates that overwrite the BN parameters by factors of
100--1000 in a single round. The model's learned normalizations are
destroyed in the first affected round; subsequent rounds cannot recover
because the same mechanism repeats.

The failure is invisible from training metrics alone: training accuracy
simultaneously reaches $85.78\%$ because workers observe their own
locally consistent BN statistics during training. Test inference fails
because the globally corrupted BN parameters normalize activations with
statistics that match no individual worker's distribution.

\textbf{Resolution.} The weighted mean server applied to shared BN
achieves $89.67\%$, marginally higher than the local BN configuration
($89.41\%$). This suggests that the collapse is caused by the adaptive
optimizer's pathological interaction with BN gradient magnitudes, not
by data heterogeneity itself. Any federated system deploying an adaptive
server optimizer on architectures with batch normalization must exclude
BN parameters from the globally synchronized payload.

\begin{table}[!t]
\caption{Ablation Study: Server Optimizer and Batch Normalization
Policy. (CIFAR-10, non-IID, NVIDIA L4. Default: FedAdam + local~BN.
$\Delta$ = change from default.)}
\label{tab:ablation}
\centering
\small
\renewcommand{\arraystretch}{1.25}
\begin{tabular}{L{3.8cm}cc}
\toprule
\textbf{Variant} & \textbf{Acc.} (\%) & \textbf{$\Delta$} (pp) \\
\midrule
Default: FedAdam + local BN      & 80.82          & --- \\
Mean server + local BN           & 89.41          & $+8.59$ \\
Mean server + Quant. + local BN  & \textbf{89.62} & $+8.80$ \\
Mean server + shared BN          & 89.67          & $+8.85$ \\
Eff.-weighted prediction         & 80.35          & $-0.47$ \\
\textbf{FedAdam + shared BN}     & \textbf{10.00} & $\mathbf{-70.82}$ \\
Aggressive quant. (2--8 bit)     & 80.04          & $-0.78$ \\
\bottomrule
\end{tabular}
\end{table}

\subsection{Multi-GPU Fast Architecture}
\label{sec:multigpu}

Table~\ref{tab:multigpu} reports results on $2\times$ T4 GPUs using Ray
actor execution with fractional GPU allocation. The current
implementation uses the same batch-level training loop and torchvision
augmentation as the single-GPU path. The reported speedups reflect Ray
scheduling and parallel worker execution rather than specialized
augmentation kernels.

\begin{table}[!t]
\caption{Multi-GPU Fast Architecture Results.
(Kaggle, $2\times$ NVIDIA Tesla T4, 20 epochs.)}
\label{tab:multigpu}
\centering
\small
\renewcommand{\arraystretch}{1.25}
\begin{tabular}{lcc}
\toprule
\textbf{Configuration} & \textbf{Acc.} (\%) & \textbf{Time} (s) \\
\midrule
\hqde{} Independent               & \textbf{94.43} & 622 \\
Shared BN + Mean server           & 87.27          & 180 \\
Shared BN + FedAdam               & 86.90          & 179 \\
Eff.-weighted prediction          & 86.00          & 179 \\
FedAvg, Mean server               & 81.82          & 179 \\
FedAvg, FedAdam                   & 81.77          & 180 \\
\bottomrule
\end{tabular}
\end{table}

\subsection{Multi-Dataset Generalization}
\label{sec:multidataset}

Table~\ref{tab:multidataset} reports \hqde{} results across five
standard datasets using 10 ensemble workers in independent mode. These
runs use the same training loop with dataset-specific dataloaders built
outside the repository; the scripts are not bundled in the current
codebase. The results suggest that accuracy drops as class count grows,
which is expected for small backbones at fixed training budgets.

\begin{table}[!t]
\caption{\hqde{} Multi-Dataset Results (10 Workers, Independent Mode).
Baseline comparison column reserved for extended version.}
\label{tab:multidataset}
\centering
\small
\renewcommand{\arraystretch}{1.25}
\setlength{\tabcolsep}{3.5pt}
\begin{tabular}{L{2.2cm}C{1.0cm}C{1.2cm}C{1.2cm}C{1.6cm}}
\toprule
\textbf{Dataset} & \textbf{Classes} & \textbf{Epochs} & \textbf{Acc.} (\%) & \textbf{Time} (s) \\
\midrule
MNIST         & 10  & 20 & \textbf{98.60} & 114.1 \\
Fashion-MNIST & 10  & 30 & 91.50          & 166.6 \\
SVHN          & 10  & 40 & 89.10          & 218.1 \\
CIFAR-10      & 10  & 40 & 81.50          & 219.4 \\
CIFAR-100     & 100 & 50 & 42.30          & 271.2 \\
\bottomrule
\end{tabular}
\end{table}

%% ============================================================
%%  VIII. DISCUSSION
%% ============================================================
\section{Discussion}
\label{sec:discuss}

\subsection{Why Hyperparameter Diversity Outperforms Seed Diversity}

The $6.86$~pp gap between \hqde{} Independent ($94.27\%$) and Deep
Ensembles ($87.41\%$) under equal compute is primarily a diversity
gap rather than a training efficiency gap. Both use 4 ResNet-18
models trained for 20 epochs on the same full CIFAR-10 dataset.

Random seed diversity in Deep Ensembles produces workers that differ
in initialization but converge toward the same broadly defined region
of the loss landscape, particularly in overparameterized networks
where multiple equivalent minima exist. Learning rate and dropout
diversity in \hqde{} produces qualitatively different optimization
trajectories: a worker with a $1.15\times$ learning rate multiplier
will overshoot low-loss basins more frequently and explore wider
regions of the landscape; a worker with a $0.10$ dropout rate will
develop denser weight representations than one with $0.15$ dropout.
The resulting ensemble members tend to specialize in different aspects
of the data distribution, and their predictions are often more
complementary than those of seed-diversified members.

\subsection{Server Optimizer Selection at Small Worker Count}

The finding that FedAdam underperforms the weighted mean by $8.59$~pp
at $K=4$ and $T=20$ is specific to the small-scale regime and should
not be generalized as evidence that adaptive server optimizers are
always inferior. The failure mechanism is sampling insufficiency:
Adam requires on the order of $1/(1-\beta_2)$ observations to
form reliable second moment estimates---with $\beta_2=0.99$, this
is approximately 100 rounds. With only 20 rounds, the moment estimates
are far from reliable, and the adaptive scaling amplifies noise.

At larger worker counts ($K>16$) and longer horizons ($T>50$),
FedAdam's moment estimates stabilize, and its benefits for
heterogeneous gradient landscapes should become apparent. A practical
guideline from these runs is to prefer the weighted mean for $K < 16$
and $T < 50$. Consider FedAdam or other adaptive optimizers only with
longer training and careful server learning rate tuning.

\subsection{The Batch Normalization Constraint as a General Warning}

The FedAdam + shared BN interaction ($80.82\% \to 10.00\%$)
is not specific to \hqde{}. Any federated learning system that:
(1) uses an adaptive server optimizer, and (2) includes batch
normalization parameters in the globally aggregated payload
will be vulnerable to this failure mode. The warning likely extends
beyond ResNet architectures to any network with batch normalization,
layer normalization with learnable parameters, or group normalization.
One mitigation is to exclude normalization parameters from the adaptive
optimizer's scope, which is a one-line configuration change in systems
with a configurable normalization policy.

This issue can be hard to spot from training metrics alone. A
practitioner monitoring only training accuracy would see $85.78\%$ and
conclude training is proceeding normally,
while test accuracy has collapsed to $10.00\%$. The diagnostic
signature is a large gap between training and test accuracy that
appears suddenly after the first few communication rounds and persists.

\subsection{Communication Compression as Regularization}

The marginal accuracy improvement from quantization (+0.21\%) over
uncompressed training, reproduced across multiple runs, is consistent
with the known regularization effects of noise injection in stochastic
optimization. Quantizing weight deltas introduces a controlled
perturbation whose magnitude is proportional to the local weight scale
(through the per-block scale factor $s_i$), not to the gradient
magnitude. This prevents large gradient components in a single round
from propagating without attenuation to the global model, effectively
smoothing the global update trajectory. At aggressive bitwidths (2--8
bit), the perturbation magnitude exceeds what error feedback can
compensate within one round, and accuracy degrades.

\subsection{Limitations}

\begin{enumerate}
  \item The primary accuracy comparison is on CIFAR-10; additional
        datasets were evaluated outside the repository.
  \item Hierarchical aggregation and Byzantine fault tolerance are
        provided as standalone modules but are not integrated into the
        core \texttt{HQDESystem} loop.
  \item Quantum aggregation utilities (superposition, entanglement,
        voting) exist in the codebase but are not benchmarked in the main
        training path.
  \item Memory comparisons depend on process layout; Ray actors run in
        separate processes while some baselines run in-process.
  \item Multi-dataset baseline comparisons are not included in the
        current code release.
\end{enumerate}

\section{Future Work}
\label{sec:future}

\textbf{Automated epoch selection.} The multi-dataset runs suggest that
optimal epoch counts vary by dataset. An early-stopping policy based on
validation trends would reduce manual tuning.

\textbf{Large-scale validation.} Scaling experiments to ImageNet and
larger backbones would test whether the ensemble and compression trends
hold at production scale.

\textbf{Integrating hierarchical aggregation.} The
\texttt{HierarchicalAggregator} module could be wired into
\texttt{HQDESystem} to evaluate tree-based aggregation at scale.

\textbf{Quantum aggregation evaluation.} The
\texttt{QuantumEnsembleAggregator} provides superposition, entanglement,
and voting variants. Benchmarking them against the efficiency-weighted
baseline remains future work.

\textbf{Byzantine fault tolerance experiments.} Controlled adversarial
injection would validate the standalone fault-tolerance module in
end-to-end training.

\textbf{Cross-domain generalization.} Applying \hqde{} to NLP and speech
would test whether the efficiency-weighted ensemble strategy generalizes
beyond vision benchmarks.

\section{Conclusion}
\label{sec:conc}

We presented \hqde{}, a distributed ensemble learning framework with
optional federated averaging and adaptive delta quantization. The
reported results indicate several practical trends.

\textbf{Hyperparameter diversity can improve ensemble accuracy.}
Independent ensemble training with learning-rate and dropout variation
reached higher CIFAR-10 accuracy than seed-only diversity in our runs.
This suggests that controlled hyperparameter variation can strengthen
ensemble complementarity.

\textbf{Adaptive delta quantization reduces communication.} Blockwise
quantization with scheduled bitwidths and error feedback reduced
communication by $7.9\times$ without accuracy loss in the reported
configuration, and by $12.6\times$ with a modest degradation. These
results suggest that the quantizer is effective in small-worker
federated settings.

\textbf{Server aggregation choice matters at small scale.} The weighted
mean outperformed FedAdam when $K=4$ and $T=20$, which suggests using the
mean when the number of rounds is limited.

\textbf{FedAdam with shared batch normalization is risky.} The reported
runs show a large accuracy drop when FedAdam is combined with shared BN.
The implementation supports local BN to avoid this interaction.

%% ============================================================
%%  REFERENCES
%% ============================================================
%% ============================================================
\balance
\bibliographystyle{IEEEtran}
\begin{thebibliography}{30}

\bibitem{lakshminarayanan2017}
B.~Lakshminarayanan, A.~Pritzel, and C.~Blundell,
``Simple and scalable predictive uncertainty estimation using deep
ensembles,''
in \emph{Proc. NeurIPS}, 2017.

\bibitem{mcmahan2017}
H.~B. McMahan, E.~Moore, D.~Ramage, S.~Hampson, and B.~A. y~Arcas,
``Communication-efficient learning of deep networks from decentralized
data,''
in \emph{Proc. AISTATS}, 2017.

\bibitem{li2020fedprox}
T.~Li, A.~K. Sahu, M.~Zaheer, M.~Sanjabi, A.~Smola, and V.~Smith,
``Federated optimization in heterogeneous networks,''
in \emph{Proc. MLSys}, 2020.

\bibitem{karimireddy2020scaffold}
S.~P. Karimireddy, S.~Kale, M.~Mohri, S.~Reddi, S.~Stich, and
A.~T. Suresh,
``{SCAFFOLD}: Stochastic controlled averaging for federated learning,''
in \emph{Proc. ICML}, 2020.

\bibitem{li2021fedbn}
X.~Li, M.~Jiang, X.~Zhang, M.~Smit, and Q.~Dou,
``{FedBN}: Federated learning on non-{IID} features via local batch
normalization,''
in \emph{Proc. ICLR}, 2021.

\bibitem{reddi2021fedadam}
S.~J. Reddi, Z.~Charles, M.~Zaheer, Z.~Garrett, K.~Rush,
J.~Kone\v{c}n\'{y}, S.~Kumar, and H.~B. McMahan,
``Adaptive federated optimization,''
in \emph{Proc. ICLR}, 2021.

\bibitem{lee2022fedndt}
S.~Lee, H.~Chung, J.~Kim, J.~Chun, J.~Moon, and T.~Hwang,
``Preservation of the global knowledge by not-true distillation in
federated learning,''
in \emph{Proc. NeurIPS}, 2022.

\bibitem{huang2017snapshot}
G.~Huang, Y.~Li, G.~Pleiss, Z.~Liu, J.~E. Hopcroft, and
K.~Q. Weinberger,
``Snapshot ensembles: Train 1, get {M} for free,''
in \emph{Proc. ICLR}, 2017.

\bibitem{zhou2022torchensemble}
T.~Chen and Z.-H. Zhou,
``{Ensemble-PyTorch}: A unified framework for ensemble learning in
{PyTorch},''
\emph{J. Open Source Softw.}, vol.~7, no.~69, 2022.

\bibitem{alistarh2017}
D.~Alistarh, D.~Grubic, J.~Li, R.~Tomioka, and M.~Vojnovic,
``{QSGD}: Communication-efficient {SGD} via gradient quantization and
encoding,''
in \emph{Proc. NeurIPS}, 2017.

\bibitem{bernstein2018}
J.~Bernstein, Y.-X. Wang, K.~Azizzadenesheli, and A.~Anandkumar,
``{signSGD}: Compressed optimisation for non-convex problems,''
in \emph{Proc. ICML}, 2018.

\bibitem{moritz2018ray}
P.~Moritz \emph{et al.},
``Ray: A distributed framework for emerging {AI} applications,''
in \emph{Proc. USENIX OSDI}, 2018.

\bibitem{blanchard2017}
P.~Blanchard, E.~M. El~Mhamdi, R.~Guerraoui, and J.~Stainer,
``Machine learning with adversaries: Byzantine tolerant gradient
descent,''
in \emph{Proc. NeurIPS}, 2017.

\bibitem{he2016resnet}
K.~He, X.~Zhang, S.~Ren, and J.~Sun,
``Deep residual learning for image recognition,''
in \emph{Proc. CVPR}, 2016.

\bibitem{krizhevsky2009}
A.~Krizhevsky,
``Learning multiple layers of features from tiny images,''
Tech. Rep., Univ. Toronto, 2009.

\bibitem{breiman2001}
L.~Breiman,
``Random forests,''
\emph{Mach. Learn.}, vol.~45, no.~1, pp.~5--32, 2001.

\bibitem{gal2016dropout}
Y.~Gal and Z.~Ghahramani,
``Dropout as a Bayesian approximation: Representing model uncertainty
in deep learning,''
in \emph{Proc. ICML}, 2016.

\bibitem{wen2020batchensemble}
Y.~Wen, D.~Tran, and J.~Ba,
``{BatchEnsemble}: An alternative approach to efficient ensemble and
lifelong learning,''
in \emph{Proc. ICLR}, 2020.

\bibitem{lin2018dgc}
Y.~Lin, S.~Han, H.~Mao, Y.~Wang, and W.~J. Dally,
``Deep gradient compression: Reducing the communication bandwidth for
distributed training,''
in \emph{Proc. ICLR}, 2018.

\bibitem{vogels2019powersgd}
T.~Vogels, S.~P. Karimireddy, and M.~Jaggi,
``{PowerSGD}: Practical low-rank gradient compression for distributed
optimization,''
in \emph{Proc. NeurIPS}, 2019.

\bibitem{karimireddy2019ef}
S.~P. Karimireddy, Q.~Rebjock, S.~Stich, and M.~Jaggi,
``Error feedback fixes {SignSGD} and other gradient compression
schemes,''
in \emph{Proc. ICML}, 2019.

\bibitem{tang2021onebitadam}
H.~Tang, S.~Gan, C.~Awan, S.~Rajbhandari, C.~Li, X.~Lian, J.~Liu,
C.~Zhang, and Y.~He,
``1-bit {Adam}: Communication efficient large-scale training with
{Adam}'s convergence speed,''
in \emph{Proc. ICML}, 2021.

\bibitem{aji2017}
A.~F. Aji and K.~Heafield,
``Sparse communication for distributed gradient descent,''
in \emph{Proc. EMNLP}, 2017.

\bibitem{li2014ps}
M.~Li, D.~G. Andersen, J.~W. Park, A.~J. Smola, A.~Ahmed,
V.~Josifovski, J.~Long, E.~J. Shekita, and B.-Y. Su,
``Scaling distributed machine learning with the parameter server,''
in \emph{Proc. USENIX OSDI}, 2014.

\bibitem{baidu2017ring}
Baidu Research,
``Bringing {HPC} techniques to deep learning,''
Tech. Blog, 2017.

\bibitem{sergeev2018}
A.~Sergeev and M.~Del~Balso,
``{Horovod}: Fast and easy distributed deep learning in {TensorFlow},''
\emph{arXiv:1802.05799}, 2018.

\bibitem{li2020ddp}
S.~Li, Y.~Zhao, R.~Varma, O.~Salpekar, P.~Noordhuis, T.~Li,
A.~Paszke, J.~Smith, B.~Vaughan, P.~Damania, and S.~Chintala,
``{PyTorch} distributed: Experiences on accelerating data parallel
training,''
\emph{Proc. VLDB Endow.}, vol.~13, no.~12, 2020.

\bibitem{dinh2020pfedme}
C.~T. Dinh, N.~H. Tran, and T.~D. Nguyen,
``Personalized federated learning with {Moreau} envelopes,''
in \emph{Proc. NeurIPS}, 2020.

\bibitem{li2021ditto}
T.~Li, S.~Hu, A.~Beirami, and V.~Smith,
``Ditto: Fair and robust federated learning through personalization,''
in \emph{Proc. ICML}, 2021.

\bibitem{micikevicius2018mixed}
P.~Micikevicius \emph{et al.},
``Mixed precision training,''
in \emph{Proc. ICLR}, 2018.

\bibitem{zaharia2010spark}
M.~Zaharia, M.~Chowdhury, M.~J. Franklin, S.~Shenker, and I.~Stoica,
``Spark: Cluster computing with working sets,''
in \emph{Proc. USENIX HotCloud}, 2010.

\end{thebibliography}

\end{document}
